\scanpage{37}
\begin{center}-- 212 --\end{center}
\hfill n$^\circ$ 391

\[
\begin{aligned}
\omega_5&=\varepsilon_4+\varepsilon_5+\varepsilon_6+\varepsilon_7+4\varepsilon_8=\\
&=8\alpha_1+12\alpha_2+16\alpha_3+24\alpha_4+20\alpha_5+15\alpha_6+10\alpha_7+5\alpha_8
\end{aligned}
\]
\[
\begin{aligned}
\omega_6&=\varepsilon_5+\varepsilon_6+\varepsilon_7+3\varepsilon_8=\\
&=6\alpha_1+9\alpha_2+12\alpha_3+18\alpha_4+15\alpha_5+12\alpha_6+8\alpha_7+4\alpha_8
\end{aligned}
\]
\[
\begin{aligned}
\omega_7&=\varepsilon_6+\varepsilon_7+2\varepsilon_8=\\
&=4\alpha_1+6\alpha_2+8\alpha_3+12\alpha_4+10\alpha_5+8\alpha_6+6\alpha_7+3\alpha_8
\end{aligned}
\]
\[
\omega_8=\varepsilon_7+\varepsilon_8
=2\alpha_1+3\alpha_2+4\alpha_3+6\alpha_4+5\alpha_5+4\alpha_6+3\alpha_7+2\alpha_8
=\tilde\alpha.
\]

(VII)\quad On sait que la demi-somme des racines positives est la somme des
poids fondamentaux, (\S\ 6, n$^\circ$ 9, prop. 8), et vaut donc :
\[
\begin{aligned}
\rho&=\varepsilon_2+2\varepsilon_3+3\varepsilon_4+4\varepsilon_5+
5\varepsilon_6+6\varepsilon_7+23\varepsilon_8=\\
&=46\alpha_1+68\alpha_2+91\alpha_3+135\alpha_4+110\alpha_5+84\alpha_6+57\alpha_7+29\alpha_8.
\end{aligned}
\]

(VIII)\quad On a vu ci-dessus que $\Gamma(R)=\Gamma_0(R)$ et l'indice de connexion
est donc 1.

\margin{\figure{\begin{tikzpicture}[scale=.5]
\draw (0,.25)--(.2,0)--(.45,.3);
\end{tikzpicture}}}
(IX)\quad La famille des exposants a 8 termes, et puisque $h=30$, les entiers
1, 7, 11, 13, 17, 19, 23, 29, étrangers à 30, doivent figurer dans cette
famille, (\S\ 8, n$^\circ$ 5, prop. 6) ; ce sont par suite tous les exposants de $W(R)$.

(X)\quad On déduit de (IX) et du \S\ 8, n$^\circ$ 3, cor. 1 de la prop. 3 que
l'ordre du groupe de Weyl $W(R)$ est :
\[
2.8.12.14.18.20.24.30=2^{14}.3^5.5^2.7.
\]

(XI)\quad Le même raisonnement qu'au n$^\circ$ 8 montre que le seul automorphisme
du schéma de Dynkin $E_8$ est l'identité, les trois chaînes issues du point de
ramification ayant des longueurs distinctes.
\mspage{36}
